% Independent audit clarification, 7 October 2026.
% For Lemma 13.26 (lem:essentiallysmoothAMMN), after the finite-range
% quasi-isogeny assertion and before the passage to circle orbits.
% This is a proposed mathematical clarification, not a formal certificate.

Here is an explicit equivariant justification for the passage to cyclic
homology. Let $C$ be the cone of the map of completed Hochschild
complexes. For every integer $N$, the preceding argument gives an
integer $a_N$ killing each of the finitely many nonzero homotopy
modules of $\tau_{\leq N}C$. The pointwise $t$-structure on
$\operatorname{Fun}(BS^1,D(\mathbf Z_p))$ gives its equivariant
Postnikov tower. Each heart object has trivial circle action:
the automorphism space of a module in the heart is discrete and
$BS^1$ is simply connected. A finite extension argument therefore
gives an integer $b_N$ such that $p^{b_N}$ annihilates
$\tau_{\leq N}C$ as an equivariant object, not merely on its
underlying homotopy groups.

Taking homotopy orbits preserves this scalar null-homotopy.
The omitted Postnikov tail remains $(N+1)$-connective after
homotopy orbits and derived $p$-completion. Indeed, homotopy
orbits of a connective spectrum are connective, and in derived
$p$-completion the potentially new bottom $\lim^1$ term is zero
because the bottom coefficient tower is the surjective tower
$\pi_{N+1}(-)/p^r$. Thus the cone of the corresponding map of
completed cyclic complexes has bounded $p$-torsion in every fixed
degree, and vanishes after inverting $p$.

Equivalently, one can take a sufficiently high finite bar skeleton
for the circle-orbit construction in each degree. Its building blocks
are finite colimits of $(S^1)^j_+\wedge X$, so the comparison commutes
with derived completion. This argument requires no uniform bound
$b_N$ as $N$ varies and no interchange of rationalization with an
infinite product. The norm fiber sequence then makes the square
between absolute and relative $HC^-\to HP$ cartesian. It does not
require that the two vertical maps of that square be equivalences
individually.

% Scope of the continuity citation:
% AMMN Prop. 4.2 applies to any qcqs scheme with bounded p-power torsion.
% The paragraph immediately before Prop. 4.10 expressly includes the
% relative Hochschild and cyclic variants. Thus this use of Prop. 4.2
% does not add a finite-presentation hypothesis on the localization B_0.
